\scanpage{61}
% manuscript line 1
près, l'opérateur usuel d'antisymétrisation dans $\bigotimes_{1\leq i\leq n}E_i$. \deleted{\unsure{C'est}}
\deleted{\unsure{facile} \illegible{1} que $a_n$ est} $a_n$ est
évidemment un \uline{projecteur continu} (car il est continu, et
\deleted{\unsure{c'est} un \unsure{proj}} \inserted{on sait que} $a_n^2=a_n$ sur le \deleted{produit} produit tensoriel algébrique
en lui-même, qui est dense), d'ailleurs hermitien
(noter que l'adjoint de $\sigma$ est $\sigma^{-1}$, à cause de
la formule $(a_{\sigma1}\otimes\cdots\otimes a_{\sigma n},b_1\otimes\cdots\otimes b_n)=(a_1\otimes\cdots\otimes a_n,b_{\sigma^{-1}1}\otimes\cdots b_{\sigma^{-1}n})$
conséquence immédiate de (1)). \deleted{\unsure{La partie}}

% manuscript line 9
\uline{Définition} Soit $E$ un espace de Hilbert, $n$ un entier $\geq0$,
on appelle \deleted{produit} \inserted{puissance} \deleted{\unsure{ext.}} extérieure hilbertienne $n$-ième de $E$,
et on note $\overset{n}{\bigwedge}\vphantom{\bigwedge}^{(2)}E$, \deleted{\unsure{la}} \inserted{l'espace de Hilbert} quotient de $\overset{n}{\otimes}\vphantom{\otimes}^{(2)}E$ par
l'opérateur $a_n$ défini par (3), mais où le produit
scalaire est multiplié par $n!$.

% manuscript line 14
Cet espace s'identifie aussi à \deleted{un} un sous-espace
de $\overset{n}{\otimes}\vphantom{\otimes}^{(2)}E$ (image du projecteur $a_n$\inserted{(au facteur $n!$ près)}). Par définition,
on voit qu'il y a une application \uline{biunivoque}
naturelle du produit extérieur algébrique $\overset{n}{\bigwedge}E$ dans
$\overset{n}{\bigwedge}\vphantom{\bigwedge}^{(2)}E$, dont l'image est identique à l'image de $\overset{n}{\otimes}E$,
donc \uline{dense} dans $\overset{n}{\bigwedge}\vphantom{\bigwedge}^{(2)}E$. On identifie donc $\overset{n}{\bigwedge}E$ à \deleted{\unsure{un}}
son sous-espace préhilbertien de $\overset{n}{\bigwedge}\vphantom{\bigwedge}^{(2)}E$, donc ce dernier
est le complété de $\overset{n}{\bigwedge}E$. Le produit scalaire \deleted{dans}
$\overset{n}{\bigwedge}E$ s'explicite aisément \deleted{\unsure{après}} sur les \deleted{\unsure{mult.}} multivecteurs
décomposables : Posons $A=a_1\otimes\cdots\otimes a_n$, $B=b_1\otimes\cdots\otimes b_n$, on a

% manuscript line 24
\[
\deleted{(a_1\wedge\cdots\wedge a_n,b_1\wedge\cdots\wedge b_n)=\left(\frac1{n!}\right)^2\sum_{\substack{\sigma\in\mathfrak S_n\\\tau\in\mathfrak S_n}}\varepsilon_\sigma\unsure{\varepsilon_\tau|\varepsilon_\sigma\,\sigma}a_1}
\]

% manuscript line 25
\[
\begin{aligned}
(a_1\wedge\cdots\wedge a_n,b_1\wedge\cdots\wedge b_n)&=n!(a_n A,a_n B)=n!(A,a_n^2B)=n!(A,a_nB)\\
&=\sum_{\sigma\in\mathfrak S_n}\varepsilon_\sigma(a_1,b_{\sigma1})\cdots(a_n,b_{\sigma n})=\det\big((a_i,b_j)\big),\quad\text{donc}
\end{aligned}
\]

% manuscript line 27
% archive sheet number: 62
